\documentclass[12pt]{amsart}


\usepackage{amssymb}
\usepackage{amsthm}
\usepackage{amsmath}
\usepackage{amsxtra}
\usepackage{mathrsfs}
\usepackage{enumerate}

\setlength{\hfuzz}{4pt}

\newenvironment{problab}[1]
{\noindent\textbf{Problem #1}.}
{\vskip 6pt}

\newenvironment{exolab}[1]
{\noindent\textbf{Exercise #1}.}
{\vskip 6pt}


\theoremstyle{remark}
\newtheorem*{soln}{Solution}

\newcommand{\C}{\mathbb C}
\newcommand{\D}{\mathfrak D}
\renewcommand{\H}{\mathfrak H}
\newcommand{\F}{\mathbb F}
\renewcommand{\P}{\mathbb P}
\newcommand{\PP}{\mathbb P}
\newcommand{\Q}{\mathbb Q}
\newcommand{\R}{\mathbb R}
\renewcommand{\SS}{\mathbb S}
\newcommand{\Z}{\mathbb Z}

\renewcommand{\Im}{\operatorname{Im}}
\renewcommand{\Re}{\operatorname{Re}}

\newcommand{\eps}{\epsilon}

\usepackage{fullpage}

\DeclareMathOperator{\Aut}{Aut}
\DeclareMathOperator{\Hom}{Hom}
\DeclareMathOperator{\impart}{Im}
\DeclareMathOperator{\repart}{Re}

\usepackage{mathpazo}

\begin{document}

\title{Belyi Maps and Dessins d'Enfants \\ Homework \#2}
\date{\today}

%\thanks{\emph{Date:} Monday, 11 February 2013}

\maketitle

\begin{exolab}{2.1}
  Show that the map
  \begin{align*}
    f: \H &\to \D\\
    z &\mapsto \frac{z-i}{z+i}
  \end{align*}
  is an isomorphism of Riemann surfaces. What is its inverse?
\end{exolab}

\begin{exolab}{2.2}
  Show that the map
  \begin{align*}
    f: \P^1 &\to \SS^2\\
    [z_0 : z_1] &\mapsto \frac{1}{|z_0|^2 + |z_1|^2} \left(2\Re(z_0 \overline{z_1}), 2 \Im(z_0 \overline{z_1}), |z_0|^2 - |z_1|^2 \right)
  \end{align*}
  is an isomorphism of Riemann surfaces.
\end{exolab}

\begin{exolab}{2.3}
  For a given $n \in \Z_{\geq 1}$, show that the map
  \begin{align*}
    f: \P^1 &\to \P^1\\
    [X:Y] &\mapsto [X^n : Y^n]
  \end{align*}
  is a morphism of Riemann surfaces.
\end{exolab}

\begin{exolab}{2.4}
  Let $E$ be the elliptic curve given by $Y^2 Z = X^3 + Z^3$. Define the map
  \begin{align*}
    f: E &\to \P^1\\
    [X:Y:Z] &\mapsto [X:Z] \, .
  \end{align*}
  \begin{enumerate}[(a)]
    \item 
      Note that this expression for $f$ does not work for the point $[0:1:0]$. How should $f([0:1:0])$ be defined in order to get a continuous map? (\textit{Hint:} Write down equations for $E$ and $f$ on the affine open set where $Y \neq 0$. Use this affine equation for $E$ to give a valid expression for $f$ on this neighborhood.)
    \item
      Using your extended definition of $f$ from the previous part, show that $f$ is a morphism of Riemann surfaces.
  \end{enumerate}
\end{exolab}

\end{document}
